\section{確率分布 : 形を言葉にし, 推論の足場をつくる}

  \subsection{問い : ヒストグラムの“形”は, 偶然か, 仕組みか}
    前章で図にしてみると, いくつかの「形」が目に入ったはずだ.  
    山がどこにあるか, どちら側に裾が長いか, ぽつんと離れた点はあるか.  
    ここで自然に湧く疑問はこうだ。

    \begin{itemize}
      \item \textbf{この形はたまたまか}. もし同じ条件で\textit{もう一度}データを集めたら, 山の位置や裾の感じは同じになるのか.
      \item \textbf{何が形を作っているのか}. 週末の増減, カテゴリーの違い, キャンペーンの有無など\textit{仕組み}が形に現れているのか, それとも小さな偶然の積み重ねなのか.
      \item \textbf{どこまでを“言える”とみなすか}. 今の図の印象を, 全体の話としてどこまで持ち出してよいのか.
    \end{itemize}

    たとえば, 20商品の購買数のヒストグラムで右に長い裾が見えたとしよう.  
    それは「人気商品の一極集中」という\textit{仕組み}が作った形かもしれないし, たまたま今週は新商品の当たり週で\textit{偶然}そう見えたのかもしれない.  
    ビン幅を少し動かすだけで山の数が変わることもあるし, サンプルが少ないと外れ値1つで印象が大きく揺れる。

    では, どう判断するか. 鍵は\textbf{取り直したときに形がどれくらい安定するか}を見ることにある.
  
    もし同じ条件でデータを何度も集められたら, そのたびのヒストグラムは少しずつ違うはずだ.  
    その「揺れ方」まで含めて形を語るために, 私たちは\textbf{確率分布}という\textit{地図}を使う.

    確率分布は, 「どの値がどれくらい出やすいか」を理想化して表す\textbf{仕組みの図解}である.  
    さらに, 同じ集め方を繰り返したとき, 平均や割合といった\textbf{統計量がどう揺れるか}（= 標本分布）もあわせて教えてくれる.  
    これが前章の「\textbf{幅をつけて話す}」の根拠になり, 次章の「\textbf{その差は偶然か}」を問う土台にもなる.

    まとめると, 私たちの問いはこう言い換えられる.
    \textit{「いま見えた形のどこまでが仕組みで, どこからが偶然か. 同じ条件で取り直したら, 形はどれくらい似るのか.」}  
    その答えに近づく道具が, この章で扱う\textbf{確率分布}と\textbf{標本分布}である.

  \subsection{目的 : 2つの分布を区別して使う}
    この章のゴールは2つ.  
    \begin{itemize}
      \item \textbf{理論分布} — 仕組みを理想化して表す「ものさし」を知り, 図の形と照らす.  
      \item \textbf{標本分布} — 同じ実験を繰り返したときに, \textbf{統計量}（平均や割合など）がどんな形で揺れるかをつかむ.
    \end{itemize}
    実データが理論分布に\textbf{ぴったり一致することはまれ}だ. それでも, 統計量の揺れ方（標本分布）は, \textbf{ある理論分布に近づく}ことが多い.  
    この事実が, 「幅で語る」「差を問う」を可能にしてくれる。
    本章は, \textbf{手元の形を見る} → \textbf{理論分布という言語で表す} → \textbf{統計量の揺れ（標本分布）を意識する} の順に進む.

\subsection{理論の最小核 : なぜ「理論分布」と「標本分布」を持ち出すのか}
  前章で見えたのは\textbf{形}だった.
  山の位置や裾の長さは, たしかに多くを語る.  
  しかし, その所見を\textbf{言葉として共有}し, なおかつ\textbf{言い過ぎずに一般化}するには, もう一歩だけ道具立てが要る.
  ここで登場するのが\textbf{理論分布}と\textbf{標本分布}である。

  はじめに強調しておく. 
  理論分布に\textbf{当てはめること自体が目的ではない}.
  私たちがやるのは, 次の順路である.
  \begin{enumerate}
    \item \textbf{まず, 手元の形を見る}. ヒストグラムや箱ひげ図から, 対称性, 裾, 多峰性, 外れの有無を素直に読む。これが\textbf{作法}（平均で語るか中央値か, 変換や外れ値処理が要るか）を決める.
    \item \textbf{次に, 形を言葉にする}. 理論分布は\textbf{モデルの形}を与える\textbf{言語}として使う.
          「釣鐘型に近い」「右に裾が長い」「回数データでばらつきがこうだ」などと短く言えるだけで, 所見は他者と\textbf{比較可能}になる.  
          ここでの役割は\textbf{説明と言語化}であり, データを型に押し込むことではない.
    \item \textbf{そして, 揺れを意識する}. 同じ条件で取り直したら, 平均や割合などの\textbf{統計量}はどのくらい\textbf{揺れる}か.  
          この「取り直したときの並び」を思い描くのが\textbf{標本分布}であり, 後章で扱う\textbf{幅づけ}や\textbf{珍しさの判断}の\textbf{物差し}になる.
  \end{enumerate}

  では, \textbf{なぜ理論分布を導入するのか}. 理由は三つに尽きる.
  \begin{itemize}
    \item \textbf{形の語彙と基準を与える}. 正規, 二項, ポアソン, 指数などは, 形を\textbf{名前で呼べる}ようにし, 観察結果を\textbf{共有可能な基準}に載せる. ぴったり一致を求めないからこそ, 語彙として有用である。
    \item \textbf{仕組みの仮説を短く描ける}. たとえば「小さな寄与の和が結果を作る」なら釣鐘型に近づきやすい, 「一定率で事象が起きる回数」ならポアソンが手がかりになる.
          これは\textbf{現象のモデル}を素描するためであり, 現実をそのまま置き換えるためではない.
    \item \textbf{使える道具と注意点が見える}. 形が分かれば, どの要約や検定が\textbf{前提に合うか}が分かり, 逆に合わないときは\textbf{変換や頑健な指標}に切り替える判断ができる.
  \end{itemize}

  そして\textbf{標本分布}を意識する理由は一つ. 私たちの推定値は\textbf{一度きりの一本の数}ではなく, 「取り直したら取り得た多くの値」の中の一点にすぎないからだ.  
  揺れの\textbf{大きさと形}を意識してはじめて, どこまで一般化してよいかを\textbf{言い過ぎずに}語れる.  

  この視点は, 次節の\textbf{正規分布と中心極限定理}で, とりわけ\textbf{平均の揺れ方が整ってくる理由}を理解する足場になる.

  \medskip
  \noindent 最小限の用語を一言でそろえる.
  \textbf{理論分布}＝現象を説明・比較するための\textbf{モデルの形}.  
  \textbf{標本分布}＝同条件で取り直したときの\textbf{統計量の揺れ方}.  
  出発点はつねに\textbf{手元の形}である.

\subsection{正規分布と中心極限定理}
    \subsubsection*{多数の小さな揺らぎの和}
      1つの測定値は, 多くの小さな要因の足し合わせで決まることが多い（製造の微差, 測定の誤差, 日々の変動など）.  
      こうした「小さな揺らぎの和」は, 釣鐘型の形を生みやすい.

    \subsubsection*{中心極限定理（芯だけ）}
      同じ集団から大きさ $n$ の標本を何度も取り直し, \textbf{平均}を毎回計算する.  
      $n$ を大きくしていくと, その平均の並び（= 標本分布）は\textbf{正規分布}に近づく.  
      元の形が多少ゆがんでいても, 観測がほぼ独立で, 条件が同程度で, 極端な外れが多くなければ, 平均は\textbf{釣鐘型}に落ち着いてくる.  
      だから本書で使った「推定値 ± およそ $2\times$ \textbf{標準誤差（SE）}」がよく効く.

    \subsubsection*{大数の法則との違い（ひと言）}
      大数の法則は「平均が真の値に近づく」話.  
      中心極限定理は「\textbf{近づき方の形}が釣鐘型になる」話. どちらも, 幅で語る根拠になる.

  \subsection{代表的な理論分布 : 仕組みから形が決まる}
    \subsubsection*{一様分布（Uniform）}
      \textbf{生成の考え方}：区間 $[a,b]$ のどの値も同じ確からしさで現れる状況. 初期化の乱数, 無作為抽選の理想像.  
      \textbf{形}：平坦. 平均は $(a+b)/2$, 広がりは $(b-a)^2/12$.  
      \textbf{目印}：ヒストグラムが山なしで平ら. 端で急に切れる.  
      \textbf{注意}：現実にぴったり平坦はまれ（丸め, 測定下限・上限の影響が出やすい）.

    \subsubsection*{二項分布（Binomial）}
      \textbf{生成の考え方}：「成功/失敗」を確率 $p$ で独立に $n$ 回繰り返し, 成功回数を数える. 例：購入に至るか, クリックするか.  
      \textbf{形}：$np(1-p)$ が大きいと釣鐘に近づき, $p$ が端に近いと片側に歪む. 平均 $np$, 広がり $np(1-p)$.  
      \textbf{目印}：回数のデータで, 0 と 1 の和が $n$.  
      \textbf{注意}：試行が独立でない, $p$ がばらつくと合わない（混合や過分散に注意）.

    \subsubsection*{ポアソン分布（Poisson）}
      \textbf{生成の考え方}：一定の平均率 $\lambda$ で, 独立にまばらに起こる出来事の\textbf{区間内の回数}. 例：到着件数, 故障件数.  
      \textbf{形}：右裾が長い. 平均と分散がともに $\lambda$.  
      \textbf{目印}：平均と分散がほぼ同じオーダー.  
      \textbf{注意}：分散が平均より大きいなら\textbf{過分散}（負の二項などを検討）.

    \subsubsection*{指数分布（Exponential）}
      \textbf{生成の考え方}：ポアソン過程で次の出来事までの\textbf{待ち時間}.  
      \textbf{形}：右に長い尾, すぐ減衰. 平均 $1/\lambda$. \textbf{記憶なし}の性質（過去の長さが次に影響しない）.  
      \textbf{目印}：生存曲線が対数軸で直線.  
      \textbf{注意}：発生率が時間で変わると合わない（ワイブルなどへ）.

    \subsubsection*{正規分布（Normal）}
      \textbf{生成の考え方}：多数の小さな寄与の和. 測定誤差や寸法のばらつき.  
      \textbf{形}：対称の釣鐘. 平均と分散で形が決まる.  
      \textbf{目印}：Q--Qプロットがおおむね直線.  
      \textbf{注意}：外れ値や重い尾に弱い（中央値や IQR 併用, 変換を検討）.

    \subsubsection*{対数正規分布（Lognormal）}
      \textbf{生成の考え方}：\textbf{掛け算の積み重ね}で大きさが決まる. 例：売上, 視聴数, 所得の一部.  
      \textbf{形}：正の値のみ, 右に長い尾. $\log$ を取ると正規に近づく.  
      \textbf{目印}：$\log$ 変換後のヒストグラムが釣鐘.  
      \textbf{注意}：平均が大きく外れ値に引っ張られる. 中央値の併記が有用.

    \medskip
    \noindent \textit{選び方の目安}：\textbf{仕組み（数え上げか, 待ち時間か, 足し算か掛け算か）}と\textbf{形（対称か, 裾の長さ）}をセットで見る.  
    ぴったり一致を目指すのではなく, \textbf{言語と基準}として使う。

  \subsection{統計量の分布 : 推論の物差し（$t$, $\chi^2$, $F$）}
    \subsubsection*{$t$ 分布}
      \textbf{場面}：平均を語るが, 全体のばらつき（母標準偏差）は未知. 標本のばらつき $s$ で代用する.  
      \textbf{形}：釣鐘に似るが\textbf{裾が太い}. 標本数が増えると正規に近づく.  
      \textbf{使いどころ}：平均の推定や平均差の比較で, \textbf{幅}や\textbf{珍しさ}を見る物差しになる.  
      \textbf{ひと言}：「自由度」はほぼ「使える独立情報の数」のこと（詳細は後章）.

    \subsubsection*{カイ二乗分布（$\chi^2$）}
      \textbf{場面}：\textbf{ばらつきの大きさ}に注目する, または「期待」と「観測」のズレを\textbf{二乗して合計}する.  
      \textbf{形}：右に偏る, 非対称.  
      \textbf{使いどころ}：分散の推定, 適合度検定や独立性の検定で, \textbf{ズレの大きさ}を測る物差しに.

    \subsubsection*{$F$ 分布}
      \textbf{場面}：\textbf{二つのばらつきの比}に注目する（分散分析, 回帰の総合的な当てはまり）.  
      \textbf{形}：右裾が長い. 上に行くほど珍しい.  
      \textbf{使いどころ}：群間の違いが\textbf{ばらつきの比}として有意かどうかを測る.

    \medskip
    \noindent \textit{要点}：$t$, $\chi^2$, $F$ は\textbf{データそのもの}の型ではなく, \textbf{統計量がどれくらいなら珍しいか}を測る\textbf{物差し}.  
    だから, データを無理にこれらに当てはめる必要はない。

  \subsection{データの分布 vs 統計量の分布 : “ぴったり”は少ないが,“揺れ方”は整う}
    \subsubsection*{データの分布（実測の形）}
      売上は右に長い（対数正規っぽい）, 待ち時間は長い裾（指数っぽい）... 現場のデータは\textbf{完全には合わない}ことが多い.  
      ゆえに, 図で確かめ, 必要なら\textbf{変換}（例: 対数）や\textbf{頑健な指標}（中央値, IQR）を使う.

    \subsubsection*{統計量の分布（標本分布）}
      いっぽう, 平均, 割合, 差, 比といった\textbf{統計量の揺れ方}は, 条件が整えば\textbf{理論分布}に落ち着く.  
      平均の揺れは\textbf{正規}に近づき, 割合や差の揺れも\textbf{式で広がりが計算}できる.  
      検定で登場する $t$, $\chi^2$, $F$ は, 「\textbf{どれくらいの値なら珍しいか}」を測る\textbf{尺}である.  
      \textit{要点}：\underline{データを無理に理論分布へ合わせる必要はない}. 必要なのは, \underline{統計量の揺れ方}をつかんで, \underline{幅と珍しさ}を判断すること.

  \subsection{最小のやり方 : 形を確かめ, 揺れを見える化する}
    \begin{itemize}
      \item \textbf{形を確認} — ヒストグラムと\textbf{密度曲線}を重ねる. Q--Qプロットで直線からのずれをざっと見る.  
            山の数, 歪み, 裾の長さを一文で要約する.
      \item \textbf{仕組みを当てる} — \textbf{数え上げ}なら二項/ポアソン, \textbf{待ち時間}なら指数, \textbf{掛け算で増える}なら対数正規を\textbf{候補語}として使う.  
            ぴったり一致は狙わず, \textbf{説明の言語}として用いる.
      \item \textbf{揺れを意識} — 平均や割合など, 使う統計量ごとに\textbf{SE}の見積もりを添え, 「推定値 ± 約 $2SE$」で\textbf{幅}を示す（根拠は CLT と $t$ の物差し）.
    \end{itemize}

  \subsection{実践 : 20商品データで体験する}
    \begin{enumerate}
      \item \textbf{購買数の形} — ヒストグラムと密度曲線. 右裾が長ければ, 対数を取って再描画. \textbf{釣鐘に近づく}か観察.
      \item \textbf{レビュー数の当てどころ} — 平均と分散を比較し, 分散が平均に近いかを確認（ポアソン的か\textbf{目安}を取る）.
      \item \textbf{標本分布の感触} — 購買数の\textbf{平均}の SE を計算し, 「推定値 ± 約 $2SE$」の\textbf{幅}を併記. 図の印象と照合.
      \item \textbf{割合でも同様に} — 星評価 $\ge 4.0$ の\textbf{割合}について, 形の所見と幅の提示を行う（端に近いと非対称に注意）.
    \end{enumerate}
    \textit{所見例}：「購買数は右裾が長い. 対数でほぼ釣鐘. 平均の SE は $\ldots$, 幅（± 約 $2SE$）を併記.」

  \subsection{実務の注意 : 正規は便利, だが万能ではない}
    \begin{itemize}
      \item \textbf{“ぴったり”を求めない} — 実測が完全に正規, はまれ. 図で形を見て, 要約や推定の道具を選ぶ.
      \item \textbf{外れ値と歪み} — 平均・分散は外れ値に弱い. 歪みが強いときは, 中央・分位や変換を併用.
      \item \textbf{独立性と混ざり} — 似た観測が連続する, 期間が混ざるなどで, 標本分布の前提が崩れる. 切り取り基準を明示.
      \item \textbf{検定の分布は“物差し”} — $t$, $\chi^2$, $F$ は物差し. データそのものの型ではない.
    \end{itemize}

  \subsection{結果と次の一手 : 分布は“形の言語”, 検定は“珍しさの言語”}
    分布を知るとは, 図の形を\textbf{言葉にする}ことだ.  
    実データの形は完璧には合わないが, \textbf{統計量の揺れ方}は理論形に\textbf{落ち着く}.  
    それが「幅で語る」を支え, 次章の\textbf{検定}で「どれくらい珍しいか」を問う土台になる.

  \subsection*{作業 : 形を見て, 言葉にして, 幅を添える}
    ECサイト20商品データで次を実施し, 小数第2位で記録せよ.
    \begin{enumerate}
      \item 購買数のヒストグラムと密度曲線. 形の要約（山の数, 歪み, 裾）を一文で.
      \item レビュー数の平均と分散を並記し, ポアソン的かどうか\textbf{目安}を記す（過分散なら注記）.
      \item 購買数の\textbf{平均}の SE を算出し, 推定値 ± 約 $2SE$ を併記.
      \item 星評価 $\ge 4.0$ の\textbf{割合}についても, 形の所見と幅の提示を行う.
    \end{enumerate}
    \textit{記録の約束}：軸・単位・期間, 切り取り基準, 欠損の扱いを明示する.

    % \section{データの傾向を捉える：記述統計から確率分布へ}

% \subsection*{問い}
% 平均や中央値, 標準偏差といった統計的要約は, 複雑なデータの集合を「ひとつの数字」に圧縮して伝えるための便利な道具である.  
% しかし, この便利さは同時に危うさもはらむ.  
% なぜなら, 要約は必然的に情報を捨てる行為だからだ.  
% 平均170cmと聞いたとき, そこに身長150cmの人と190cmの人が混ざっている可能性や, ほとんど全員が169〜171cmの範囲に収まっている可能性は区別できない.  

% この章で私たちが探るのは, **「平均や中央値の背後に潜む, データ全体の姿」**である.  
% それを理解するための武器が「確率分布」である.

% \subsection*{素朴な方法と限界}
% 最初の一歩として有効なのはヒストグラムだ.  
% データを区間に分け, その中に含まれるデータ数を棒の高さで表せば, 山や谷, 左右の広がりが視覚的に掴める.  
% しかし, ヒストグラムには次のような限界がある.  

% \begin{itemize}
%   \item 区間幅（ビン幅）を変えると形が大きく変わる.
%   \item データが少ないと偶然の変動に形が左右されやすい.
%   \item あくまで現在あるサンプルに基づく形であり, 無限にデータを集めたときの「理想的な姿」はわからない.
% \end{itemize}

% たとえば, サンプル数が30のときと3000のときでは, 同じ現象を測ってもヒストグラムの滑らかさや形の安定性はまるで違う.  
% そこで登場するのが, **「理論的な分布」**という発想である.

% \subsection*{理論の最小核}
% 理論分布とは, データが取り得る値とその確率の対応を, 数学的な関数として記述したものだ.  
% 離散データの場合は確率質量関数, 連続データの場合は確率密度関数で表す.  
% 共通する条件は, その合計（または面積）が必ず1になることだ.  

% 重要なのは, 理論分布は現実をそのまま描写したものではないという点だ.  
% むしろ, 実験や観測の背後にある「理想化された確率的な仕組み」を描き出す道具である.  
% 現実のコインは完全に対称ではないし, 投げ方や環境によって微妙に確率が変わるかもしれない.  
% それでも「表が出る確率は1/2」という単純なモデルを仮定することで, 複雑な現象を数理的に扱えるようになる.  

% この物差しをもとに, 実際のデータがどの分布に近いのかを調べれば, その現象の性質や適用可能な統計手法が見えてくる.

% \subsection*{正規分布と中心極限定理}
% 多くの理論分布の中でも, 統計学で特別な地位を占めるのが**正規分布**だ.  

% \subsubsection*{正規分布の特徴}
% 正規分布$N(\mu, \sigma^2)$は, 平均$\mu$と分散$\sigma^2$という2つの数値だけで完全に形が決まる連続分布である.  
% 形は左右対称の釣鐘型で, 中心から離れるほど値の出現確率が急速に下がる.  
% 平均が変われば山の位置が水平移動し, 分散が変われば山の幅（広がり）が変わる.  
% 確率密度関数は次の式で表される.

% \[
% f(x) = \frac{1}{\sqrt{2\pi\sigma^2}} \exp\left( - \frac{(x-\mu)^2}{2\sigma^2} \right)
% \]

% この式に含まれる$\exp(-x^2)$という形は偶然ではない.  
% 左右対称で, 加法に対して形が保たれ, しかも分散が有限な分布を求めると, この形以外には存在しない.  
% また, 平均と分散が与えられたときにエントロピー（不確かさ）が最大になる分布もこの形である.

% \subsubsection*{なぜ頻繁に現れるのか}
% 答えは**中心極限定理**にある.  

% \begin{thm}
% 平均$\mu$, 分散$\sigma^2$をもつ同一分布から独立に$n$個のデータを抽出し, その平均をとる.  
% $n$が大きくなるにつれて, その平均の分布は正規分布$N(\mu, \sigma^2/n)$に近づく.
% \end{thm}

% つまり, 元の分布の形がどんなに歪んでいても, 平均を取れば正規分布に近づく.  
% これが, 正規分布が自然界・社会・工業製品のあらゆる場面で出現する理由だ.  

% \subsubsection*{直感的な説明}
% 身長の例を考えよう.  
% 1人の身長は遺伝, 栄養, 運動習慣, 健康状態など数え切れない要因が少しずつ影響して決まる.  
% それぞれの要因は大きな偏りをもたらすわけではなく, 小さな変動を加えるだけだ.  
% この「多数の小さな揺らぎの和」が, 釣鐘型の正規分布を生み出す.  

% \subsubsection*{計算実験の例}
% 次のシミュレーションは中心極限定理の威力を示す.  
% \begin{enumerate}
%   \item $0$から$1$までの一様乱数を12個発生させ, それらの和を求める.
%   \item この操作を1万回繰り返す.
%   \item 結果をヒストグラムにすると, 元の乱数は平坦だったのに, 和は滑らかな釣鐘型になる.
% \end{enumerate}

% \subsection*{誤差論との関係}
% 正規分布は, 「誤差をどう扱うか」という古典的な問題とも深く結びついている.  
% 測定値は必ず何らかの誤差を含むが, その多くは小さな要因がランダムに重なった結果である.  
% このとき, 誤差の分布は理論的には正規分布になる――これはガウス以来の誤差論の基本原理だ.  

% 実務では, 測定値の誤差が正規分布であると仮定すると, 平均値や分散の推定式が非常にシンプルになり, 検定や区間推定も効率的になる.  
% ただし, 実際の誤差が「正規的な性格」を持っていない場合もあり, たとえばセンサーの異常や人間の記録ミスのような外れ値が混ざると, 平均値は大きく引きずられてしまう.  
% こうした場合はロバスト統計の出番であり, 「正規分布はありがたいけれど, 盲信は禁物」というわけだ.

% \subsection*{実務上の注意点（少しだけユーモラスに）}
% 正規分布はあまりに有名すぎて, ときどき「統計＝正規分布」と勘違いされることがある.  
% しかし, 実際のデータの多くは完全な正規分布から微妙に（あるいは大胆に）外れている.  
% 待ち時間のデータは長い尻尾を引く指数型だったり, 売上額は右にびよーんと伸びた対数正規型だったりする.  

% 現場ではこんな会話がよくある.  
% 「このデータ, 正規分布に従っていることにしましょうか？」  
% ――これは魔法の呪文ではなく, 単なる作業上の仮定だ.  
% もちろん仮定は便利だが, 無理やり当てはめると後で痛い目を見る（検定結果が使い物にならなくなる, など）.  

% したがって, 正規分布を仮定するときは, かならずヒストグラムやQ-Qプロットで形を確認し, 必要に応じて歪度・尖度なども見ておくこと.  
% 正規分布は「万能の鍵」ではなく, あくまで「鍵束の中の1本」だという意識を持っておくと, データ分析の事故はぐっと減る.

% \subsection*{結果と次の一手}
% 分布の形を知ることは, 単なる分類ではない.  
% それは, 適用できる統計手法, 推定の精度, 検定の前提条件を判断するための基礎である.  
% 次章では, この分布を物差しにして「その差は偶然か必然か」を見極める統計的検定の世界に進む.



% \section{データの傾向を捉える：記述統計から確率分布へ}

% 平均値, 中央値, 標準偏差などの指標は, 複雑なデータをひとつの数字にまとめてわかりやすく表現してくれる. 
% これは, たとえば初めて行ったケーキ屋さんで何を買うか迷ったときに, その店の売れ筋ランキングを参考にするようなものだ. 
% 一方で, ケーキにもそれぞれ異なった魅力があるように, 実際のデータにもより細かな傾向や性質がある. 
% それらを本当に理解しようとすれば, 平均などの「数字の要約」だけでは足りないことも多い． 
% 前章ではデータを「可視化」することで, より全体的な傾向や特徴を視覚的に掴もうと試みた. 
% この章ではそこから一歩進み, データのばらつき方やパターンを理論的に捉える方法について述べる．  
% そのために, 「データの分布」について考えてみよう．

% \subsection{ヒストグラムと確率分布}

% データの分布を視覚的にとらえる方法として，最も基本的なのがヒストグラムである．  
% ヒストグラムは，データの値をいくつかの区間に分けて，それぞれの区間にいくつのデータが含まれているかを棒グラフのように表したものだった．
% このヒストグラムの形に注目して，「もしデータが大量にあったら，どんな形になるだろう？」と考えると，そこから「確率分布」という考え方につながっていく．  
% たとえば連続的なデータの場合，ヒストグラムの各棒を細かくして，面積の合計が$1$になるように調整すれば，それは「確率密度関数」と呼ばれるものに近づくのであった．  
% これは，「ある値の近くでデータが現れやすいかどうか」を表すものである．
% ある値までの面積を計算することによって, たとえば身長$160$cmまでの人がどのくらいの「割合」いるか, といったことがわかる.
% 一方で，サイコロの出目のような離散的なデータでは，各値に対応する確率をそのまま棒で表した「確率質量関数」が使われる．  
% これらを合わせて「確率分布」と呼ぶ．

% 理論的な分布は，ある種の理想的な仮定のもとで考えられている．  
% たとえば「コインを投げたときに表が出る確率はちょうど$1/2$」というような仮定がある．  
% 現実ではコインの投げ方や場所など，さまざまな要因が結果に影響するかもしれないが，そのようなノイズを無視して考えるのが理論分布である．

% そして，もしこのような理論的な試行を何度も繰り返したとすれば，得られるデータの分布は理論分布に近づいていく．  
% これは「大数の法則」によって保証されている．
 
% ここでひとつ用語についての注意をしておこう.
% これからときどき, なんらかのデータが「ある理論的な分布に従う」ということがあるが，それは「そのデータの分布が, ある理論的な分布と同じような形をしている」という意味である．


% \subsubsection*{理論分布を使うときの心構え}
% これから代表的な理論分布を紹介するが，そのときにひとつだけ注意しておいてほしいことがある．  
% それは，「すべてのデータが，どれかの理論分布に必ず従っている」というわけではない，ということである．

% 上記のような「確率分布」についての導入を読むと, あたかも全てのデータに対して, そのデータの背景には何らかの理論的な分布が潜んでいるように感じる.
% そして, データについて知りたければその分布を調べればよいのであり, データをどんな分布に当てはめるかが中心的問題なのでは？と思ってしまうことがある.

% たしかに，データがある理論分布に近い形をしていることはある．  
% しかし，実際の世界はもっと複雑である.
% データがなんらかの理論分布にいつも当てはまるとは限らないし，そもそも「どの分布に従っているか」は完全にはわからないことも多い．  
% また，統計的に何かを判断したいときに，データがどの理論分布にぴったり当てはまるかを調べる必要があるとは限らない．

% この章で紹介する分布は，「よく知られた形の代表例」であり，現実のデータを無理にそれに当てはめるためのものではない．  
% それよりも，データを眺めるときの基準として役立ったり，統計的な検定を行うときの前提として使われたりするものだと考えてほしい．

% とはいえ, もちろん自分のデータがどの分布に近いのかを調べることにも重要な意味をもつのは確かであり, それが気になるのも人情というものである.  
% この章の後半では，そのための方法についても簡単に紹介する． 

% \subsection{確率分布}
% 確率分布には用途によっておおまかに「データを表す分布」と「統計的検定のための分布」に分けられる. 
% これらの種類分けは一般的なものではないが, 種類ごとに紹介した方がわかりやすいような気がするのでそのように書いてある.
% 最初に紹介する正規分布はデータを表すことにも, 検定のためにも重要なので特別扱いをしている. 

% \subsubsection{正規分布}

% 正規分布（いわゆる「釣鐘型」の形をした分布）は, 統計学において最も重要な分布である. 
% はじめて統計を学ぶときは, この分布について「だいたい分かった気がする」だけで十分といえるくらいだ. 
% 実際，正規分布は現実のさまざまな現象において自然に現れる分布であり,  
% 人間の身長, 雨粒の大きさ, 製品の重さなど, 多くの測定値がこの形に近い分布を示すと言われている．

% ではなぜ，このような分布がよく現れるのだろうか？  
% その答えのひとつが，「中心極限定理」にある．

% \begin{thm}
%   平均 $\mu$，分散 $\sigma^2$ をもつ同じ分布から，独立にデータをたくさん（$n$個）取り出して，その平均をとる．
%   このとき，$n$ が大きくなればなるほど，その平均の分布は正規分布に近づいていく．
% \end{thm}

% この定理が教えてくれるのは，いろいろな要因が重なって生まれたデータの平均値は，たとえ元のデータが正規分布に従っていなくても，
% その平均を何度も取り出して集めていくと最終的には正規分布に近づいていくということだ．
% （データをたくさん集めるのではない. 「データの平均」をたくさん集めるのだ．）  
% この結果は一見不思議に思えるかもしれないが，実はとても自然なことでもある．
% なぜなら，私たちが現実世界で測る多くのものは，「ひとつの原因」ではなく，無数の小さな要因が重なった結果だからだ．

% たとえば，ある製品の重さを測ったとしよう．
% その重さには，材料のわずかなばらつき，加工のずれ，測定器の誤差，環境条件など，数えきれないほど多くの要因が影響している．
% これらがバラバラな方向に，ちょっとずつ作用すると，その誤差の合計として現れる最終的な重さは「全体の平均」の近くに集まりやすく，
% 少し重かったり軽かったりする製品が，同じくらいの割合で現れる様子はイメージしやすいのではないだろうか．
% このような要因の重なりが, 結果として釣鐘型の正規分布を自然に形作るのである．

% 中心極限定理は，こうした「多くの小さな揺らぎの合計が生み出す形」として，正規分布が非常に基本的な存在であることを数学的に裏づけている．
% だからこそ，正規分布はあらゆる場所に現れるのだ．
% それは人間が自然を観測し，測定し，理解しようとするかぎり，避けがたく現れる形なのである．
% 実践的にも, データの平均を取って, その分布をみることは意外とある. 
% たとえば, 各人のデータを一定期間で平均を取って, それをいろんな人のグループで比較するなどというときはこの分布が活躍する. 

% ただその一方で，正規分布がすべてのデータに当てはまるわけではないことには注意が必要だ．  
% 「なにかの大きさを測る」のとは少し違うこと, たとえばなんらかの出来事が起きた回数（事故の件数）などや，
% 次の出来事が起きるまでの時間（待ち時間など）といった現象では，正規分布ではなく他の分布が使われることもある．

% このように，正規分布は実際のデータを表すうえで非常に役立つが，それだけではない．
% 正規分布は，統計的検定の多くの場面でも，基準や前提として用いられる．
% たとえば，ある集団の平均と仮説上の値が異なるかどうかを判断するとき，その背後では「標本平均は正規分布に近い」という性質が使われている．
% つまり正規分布は，単にデータの傾向を表すだけでなく，「判断の基準としての分布」という側面も持っているのである．

% 正規分布がもうひとつ特別である理由は，「平均と分散という2つの数値だけで形が完全に決まる」という点にある．  
% 平均が変わると分布の中心が動き，分散が変わると広がり方が変わるが，それ以外のかたちは変わらない．  
% この性質は，データが正規分布に従うと仮定できる場合に，ごく少ない情報からでも有益な推測ができることを意味している．

% ただし，平均と分散だけで正規分布の全てがわかるというわけではない. 
% 正規分布はそれ以上に確率分布としての豊かな情報を持っており, 推定や検定など，多くの統計的手法が正規分布を前提として構築されている．
% 平均と分散によって一意に決まるというシンプルな性質が，かえって曖昧になりがちなデータ分析を明確に捉えなおすための土台となっている．
% これは，住所や電話番号といった個人情報で誰かを特定できても，その人の性格や考え方までは分からないのと似ている．

% このように，正規分布は実データの傾向を理解するためにも，統計的な判断の基準としても，中心的な役割を果たしている．

% \begin{quote} \textbf{コラム：なぜ正規分布は $e^{-x^2}$ なのか？}
%   正規分布は，$e^{-x^2}$ という少し不思議な形をした関数によって表されている．
%   なぜ，この関数が選ばれているのだろうか？
  
%   正規分布が現れる理由は，誤差などの小さなばらつきをたくさん足し合わせた結果として自然に生じることにあった．
%   そのような積み重ねが繰り返されるとき，足し合わせに対して分布の形が保たれていてほしいという条件はごく自然に要請される．
%   そしてこの条件を満たし, 分散がきちんと定義される「左右対称」な関数は，実は $e^{-x^2}$ の形をしたものしかない．
%   さらに，自然さという意味では, 「平均と分散が決まっているときに，それ以外に一切の恣意性を含まない分布」を求めたときにも，やはりこの関数が現れる．
%   これはエントロピーという考え方で導かれるものである. 
  
%   $e^{-x^2}$ という関数で定義される分布が「正規」な分布して有用なのは，ばらつきを扱う上で最も自然で恣意性のない形だからなのである． 
% \end{quote}

% \subsubsection{データ生成を表す分布}

% ここでは, 実際によく見られる現象を表すために, 統計学でよく使われる代表的な確率分布を紹介する．
% 繰り返しになるが, 現実のデータが必ずしもこれらの分布のいずれかにぴったり当てはまるわけではない．
% それでも, これらの分布の特徴を知っておくと, データの性質を理解したり, 観察されたデータがどのような仕組みから生まれたのかを考える手がかりになる．

% \begin{itemize}
%   \item \textbf{ポアソン分布}：  
%   一定の時間や空間の中で, ある出来事が何回起きたかを扱う分布． 
%   単位時間あたりに平均$\lambda$回起きると仮定すると，そのときの回数のばらつきを表す． 
%   例：1時間あたりの来客数，1日あたりの交通事故の件数など．  
%   《ここに図を入れる：ポアソン分布の棒グラフ（平均値の違いによる例など）》

%   \item \textbf{二項分布}：  
%   成功か失敗かといった, 二つの結果のうちどちらかが起きる試行を何度か繰り返したとき，成功の回数がどのように分布するかを表す．  
%   例：10回コインを投げて表が出る回数など．  
%   《ここに図を入れる：$n = 10, p = 0.5$ の二項分布の棒グラフ》

%   \item \textbf{指数分布}：  
%   次の出来事が起こるまでにかかる時間に関する分布．
%   単位時間あたりに平均$\lambda$回起きるなら，次が起きるまでの時間は指数分布に従う．
%   一定時間内の回数を表すポアソン分布と「表裏」の関係にあると言える．
%   例：次の顧客が来店するまでの待ち時間など．   
%   《ここに図を入れる：指数分布の滑らかなグラフ（$\lambda$ の違いによる例）》

%   \item \textbf{一様分布}：  
%   どの値も等しい確率で現れる分布．
%   なんらかの数値を初期設定する際に用いられることもある. 
%   例：理想的なサイコロの出目（1から6までが全て同じ確率で出る）など．  
%   《ここに図を入れる：離散型と連続型の一様分布の例（棒グラフ・長方形）》

%   \item \textbf{ベータ分布}：  
%   確率や割合のように，0から1の範囲に収まる量を扱う分布．  
%   観測回数が少ないときの不確かさを表すのに使われる．  
%   パラメータの値に応じて，山型，U型，平坦などさまざまな形になる． 
%   《ここに図を入れる：ベータ分布（パラメータによって形が変わる例を2, 3）》

%   \item \textbf{対数正規分布}：  
%   データの対数（log）をとると正規分布に従うような分布．  
%   所得，再生回数，販売数など，右に長く尾を引く「非対称な分布」にしばしば見られる．  
%   小さい値が多く，大きい値はまれに出現するという性質を持つ． 
%   《ここに図を入れる：正規分布との比較つきの対数正規分布の例》
% \end{itemize}


% \subsubsection{統計量の分布}

% 上記で紹介した分布は, 実際のデータがどのように分布しているかを見るためのものだった．
% 一方で, 統計的検定を行うときには，データそのものではなく，そこから計算される「統計量」がどのような分布に従うかを考える必要がある．
% ここでいう統計量とは, 後の章で詳しく扱うが, 平均や分散など，データ全体を要約するために計算される値のことである．
% 同じような調査や実験を何度も行ったとしたら，統計量の値も毎回少しずつ変わるだろう． 
% 統計的な判断を下すときには，「この統計量はどのくらいの確率でこの値になるのか？」という発想が必要であり, 
% そのときに検定や推定の正当性の根拠として使われるのが，これから紹介する「統計量の分布」である．

% \begin{itemize}
%   \item \textbf{$t$分布}：  
%   標本のサイズがそれほど大きくないときに，偶然で片づけられないほど平均の差に意味があるかどうかを検定する際に使う分布．  
%   正規分布とよく似た形をしているが，裾がやや広く，小さなデータでは不確かさが大きいことを反映している．
%   標本数を増やすと正規分布に近づいていく. 
%   \textit{（例：t検定に用いられる）}

%   \item \textbf{カイ二乗分布}：  
%   データのばらつき具合や，ある理論分布との適合の程度を評価するときに使う分布．  
%   たとえば，観測された回数と期待される回数のずれが「偶然にしては大きすぎるか」を判断するときに用いられる．．  
%   \textit{（例：適合度検定・独立性の検定など）}

%   \item \textbf{$F$分布}：  
%   複数のグループのばらつき（分散）に注目し，それらを比較して統計的に意味があるほど異なるかどうかを検討する際に使う分布．  
%   \textit{（例：主に分散分析（ANOVA）で用いられる）}
% \end{itemize}

% これらの分布は，いずれも「統計量がどの程度の値をとるのが自然か（それとも不自然か）」を判断する基準として使われる．  
% 正規分布が「データそのもののばらつき」を表すとすれば，$t$分布やカイ二乗分布，$F$分布は「判断のための物差し」を提供してくれる分布だといえる．



% \subsection{現実のデータはどの分布に近いのか}

% ここまでに紹介してきた理論的な分布は, あくまで「理想的な形」であり，実際のデータがそれとまったく同じになることはほとんどない．
% それでも，データの全体的な傾向をつかむために，「どの分布におおまかに近いか」を確認しておくことには意味がある．  
% ぴったり一致する必要はなく，「大まかに似ている」ことがわかれば，それに基づいて適切な検定方法を選んだり，予測や推定の方針を立てたりすることができる．

% 自分のデータがどの分布に近いかを判断するための代表的な方法には，以下のようなものがある：

% \begin{itemize}
%   \item \textbf{ヒストグラムによる視覚的判断}：  
%   データの形を棒グラフとして描き，山の数や対称性，ばらつきの広がりを観察する．  
%   たとえば，正規分布に近いデータは釣鐘型になることが多い, などである. 
%   眼で見て判断するというきわめて単純な方法であるが意外と役に立つ.
%   全然違う分布を当てはめてしまったときなどは一目瞭然である.

%   \item \textbf{累積分布関数（CDF）を用いた確認}：  
%   「ある値以下になる割合」をグラフとして描く．   
%   理論分布の累積分布関数のグラフと並べて比較することで，全体の形がどのくらい似ているかを視覚的に比較できる．

%   \item \textbf{Q-Qプロット（Quantile-Quantile Plot）による正規性の確認}：  
%   データの分位点（何パーセント地点か）と理論分布の分位点を対応させてグラフを描く方法．  
%   点がおおむねひとつの直線上に並べば，その理論分布に近いと判断できる．  
%   （特に正規分布との比較によく用いられる）


%   \item \textbf{適合度検定，歪度（Skewness），尖度（Kurtosis）による数値的判断}：  
%   グラフだけでは判断が難しい場合に，データが特定の分布からどのくらい離れているかを数値的に評価する方法．  
%   歪度（左右の偏り）や尖度（尖り具合）などの統計量を使って，より客観的に分布の形を評価する．  
%   （適合度検定についての詳細は次章で扱う）
% \end{itemize}

% どの方法にも一長一短があるため，ひとつの方法に頼るのではなく，視覚的・数値的な観点の両方から判断するのが望ましい．
% 大切なのは，「どの分布にぴったり当てはまるか」ではなく，「どの分布におおまかに近いか」という感覚をもつことである．


% \textbf{メモ}\\
% ヒストグラムをつくる例を入れると良いかもしれない.簡単なデータ, たとえば所得のデータがいいかも.所得の分布が指数分布になっているとか.
% 連続分布と離散分布の違いについて触れたほうがいいかもしれない.
% 最後の分布の当てはめについては, イェンゼン-シャノン距離などについても触れたほうがよいか？